% Source-aligned English for printed p.169.
\SourcePage{169}

\LayerTarget{tohoku:proof:R-quotient-ZYm:p2}
\noindent zero. On the other hand, the restriction of $R_m/R_{m-1}$ to
$Y_m=U_m\cap \mathrm C V_{m-1}$ is isomorphic to that of $R_m$ (since that
of $R_{m-1}$ is zero) and is generated there by the restriction of the section
$f_m$; hence it is isomorphic to the constant sheaf $Z$.

\LayerTarget{tohoku:proof:3.6.1:p3}
This proves that $R_m/R_{m-1}$ is isomorphic to $Z_{Y_m}$ and reduces us to
proving that $T^{p+1}(Z_Y)\in C''$ when $Y$ is a locally closed subset of
$X$. We may write $Y=V\cap \mathrm C U$, where $U$ and $V$ are open subsets
of $X$. Let $W$ be the open subset $U\cup V$ of $X$. There is an evident exact
sequence
\[
0\longrightarrow Z_U\longrightarrow Z_W\longrightarrow Z_Y
\longrightarrow0,
\]
and hence an exact sequence
\[
T^{p+1}(Z_W)\longrightarrow T^{p+1}(Z_Y)
\longrightarrow T^{p+2}(Z_U).
\]
By hypothesis the two outer terms of this exact sequence belong to $C''$;
therefore so does $T^{p+1}(Z_Y)$. This completes the proof of the proposition.

\LayerTarget{tohoku:proposition:3.6.2}
\noindent\textsc{Proposition 3.6.2.} --- \textit{Let $C$ and $C'$ be two
abelian categories satisfying axiom AB 5) (cf. 1.5), and suppose that $C$
admits a generator (cf. 1.9). Let $T$ be a covariant functor from $C$ to $C'$.
For the right derived functors $R^pT$ to commute with inductive limits, it is
enough that:}
\begin{description}
\LayerTarget{tohoku:subitem:3.6.2:a}
\item[\textit{a)}] \textit{$T$ commute with inductive limits;}
\LayerTarget{tohoku:subitem:3.6.2:b}
\item[\textit{b)}] \textit{if $M=\varinjlim_i M_i$ in $C$, where the $M_i$
are injective, then $M$ is $T$-acyclic, i.\,e.
$R^pT(M)=0$ for $p>0$.}
\end{description}

\begin{quote}
\small
\textbf{Editorial qualification.} With the definition on p.~143, for a
covariant functor that is assumed only additive, the canonical morphism
$T(M)\longrightarrow R^0T(M)$ need not be an isomorphism. Thus the
vanishing $R^qT(M)=0$ for $q>0$ required in b) does not by itself supply
the degree-zero comparison needed for the equality
$R^pT(A)=H^p(T(C))$ used below. The proposition and its proof are valid
if one adds the hypothesis that $T$ is left exact. Without this
hypothesis, b) must instead be strengthened by requiring, for every
$M=\varinjlim_i M_i$ occurring there, both that
$T(M)\longrightarrow R^0T(M)$ be an isomorphism and that
$R^qT(M)=0$ for $q>0$. This qualification changes neither a), nor the
compatible-resolution construction, nor the resulting comparison
isomorphisms under either qualified hypothesis.
\end{quote}

\LayerTarget{tohoku:proof:3.6.2:p1}
(Condition b) is plainly necessary for the conclusion to hold, and so is a)
when $T$ is left exact, hence $R^0T=T$.) First note that the $R^pT$ are
well-defined by Theorem 1.10.1. Let $(A_i)_{i\in I}$ be an inductive system in
$C$, and let $A$ be its inductive limit. We wish to prove that the natural
morphisms

\LayerTarget{tohoku:eq:derived-colimit-comparison}
\[
\varinjlim_i R^pT(A_i)\longrightarrow R^pT(A)
\]
are isomorphisms.

\LayerTarget{tohoku:construction:inductive-injective-resolutions}
We first show that there exists an inductive system $(C_i)_{i\in I}$ of
complexes with values in $C$, together with an ``augmentation''
$(A_i)\longrightarrow(C_i)$, such that for every $i\in I$ the map
$A_i\longrightarrow C_i$ is a right resolution of $A_i$ by injective objects.

\LayerTarget{tohoku:construction:diagram-category-resolution}
Indeed, consider the category $I(C)$ of inductive systems indexed by $I$ with
values in $C$. It is a diagram category and, by Proposition 1.6.1, satisfies
the same hypotheses as $C$. By Theorem 1.10.1, every object of this category
therefore admits a right resolution by injective objects. Moreover, if $(M_i)$
is an injective object of $I(C)$, then the $M_i$ are injective objects of $C$
(this is a general fact for diagram categories indexed by a scheme $\Sigma$
satisfying the general conditions of Proposition 1.9.2).

\begin{quote}
\small
\textbf{Editorial correction.} The appeal to Proposition 1.9.2 is not valid
as printed: its assertion about the stated family of generators is false in
general. For $i$ in the small filtered poset $I$, let
$F_i:C\longrightarrow I(C)$ be the left adjoint to evaluation at $i$, defined
by
\[
F_i(A)_j=
\begin{cases}
A,&i\leq j,\\
0,&\text{otherwise,}
\end{cases}
\]
with identity transition maps wherever both values are $A$. This functor is
exact objectwise. If $E$ is a generator of $C$, then
$\bigoplus_{i\in I}F_i(E)$ is a generator of $I(C)$. Since each $F_i$ is an
exact left adjoint, its right adjoint, evaluation at $i$, sends injectives of
$I(C)$ to injectives of $C$. Thus $I(C)$ has the generator required by
Theorem 1.10.1, evaluation preserves injectives, and the remaining argument is
valid.
\end{quote}

Take a right resolution of $(A_i)$ by injective objects of $I(C)$,

\LayerTarget{tohoku:eq:diagram-injective-resolution}
\[
0\longrightarrow(A_i)\longrightarrow(C_i^0)\longrightarrow(C_i^1)
\longrightarrow\cdots.
\]

\LayerTarget{tohoku:eq:component-injective-complex}
For each $i$, let $C_i$ be the complex
\[
0\longrightarrow C_i^0\longrightarrow C_i^1\longrightarrow\cdots.
\]
Then the system $(C_i)$ has the required property.

\LayerTarget{tohoku:eq:colimit-resolution}
Since the functor $\varinjlim$ on $I(C)$ is exact (Proposition 1.8), we obtain
a resolution $C$ of $A$ by the objects
\[
C^p=\varinjlim_i C_i^p.
\]

\LayerTarget{tohoku:eq:T-acyclic-computation}
By condition b), the $C^p$ are $T$-acyclic, and therefore
\[
R^pT(A)=H^p(T(C)).
\]

\LayerTarget{tohoku:eq:T-colimit-complex}
Since $T$ commutes with inductive limits by a),
\[
T(C)=\varinjlim_i T(C_i).
\]
The exactness of $\varinjlim$ on $I(C')$, which follows from axiom AB 5) by
Proposition 1.8, now gives

\LayerTarget{tohoku:eq:cohomology-colimit-comparison}
\[
H^p(T(C))=\varinjlim_i H^pT(C_i).
\]
But since $C_i$ is the complex

\TohokuContinuation{170}
